\section{Energy facts for adaptive relaxation}
\label{sec:energy}

The next statements apply to the unregularized quadratic $f$ in the shared
model.  They do not assume monotone residual signs.
Define the error energy
\begin{equation}
 \mathcal E(\widehat x)
 :=f(\widehat x)-f(x^0)
 =\frac12(\widehat x-x^0)^\mathsf TQ(\widehat x-x^0).
 \label{eq:error-energy}
\end{equation}

\begin{lemma}[Exact coordinate-energy ledger]
\label{lem:energy}
Let $\widehat x$ be any iterate, $r=b-Q\widehat x$, and let
$\widehat x^+$ be the row update \eqref{eq:row-update} on $u$.  For every
$0<\omega<2$,
\begin{equation}
 f(\widehat x)-f(\widehat x^+)
 =\frac{\omega(2-\omega)}{2Q_{uu}}r_u^2
 =\frac{\omega(2-\omega)}{1+\alpha}r_u^2.
 \label{eq:energy-drop}
\end{equation}
Consequently, along any finite adaptive run,
\begin{equation}
 f(\widehat x_0)-f(\widehat x_T)
 =\sum_{t=0}^{T-1}
 \frac{\omega_t(2-\omega_t)}{1+\alpha}r_{t,u_t}^2.
 \label{eq:telescope}
\end{equation}
The right-hand side is computable from the live local residuals and does not
require knowing $x^0$.
\end{lemma}

\begin{proof}
Set $z=\omega r_u/Q_{uu}$.  The exact quadratic expansion gives
\[
 f(\widehat x+ze_u)-f(\widehat x)
 =z\nabla_u f(\widehat x)+\tfrac12Q_{uu}z^2
 =-\frac{\omega(2-\omega)}{2Q_{uu}}r_u^2.
\]
The shared matrix has $Q_{uu}=(1+\alpha)/2$.  Summing proves the telescoping
identity.
\end{proof}

\begin{corollary}[Myopic descent selects the exact rung]
\label{cor:myopic}
At a fixed iterate and coordinate, immediate objective decrease per charge
is uniquely maximized over $0<\omega<2$ by $\omega=1$.
\end{corollary}

\begin{proof}
The residual and charge are fixed, while
$\omega(2-\omega)=1-(\omega-1)^2$ is uniquely maximized at one.
\end{proof}

Thus no controller based only on the current update's energy efficiency can
explain the measured advantage of over-relaxation.  The benefit must appear
over a block: over-relaxation changes which coordinates become live, how far
the frontier travels, and which expensive rows are revisited.  Equation
\eqref{eq:telescope} is still useful as an auditable block-progress
denominator, but comparing modes from the same starting state requires a
probe, a branch, or a structural theorem.
